\origpage{327}{N. Robertson, P.D. Seymour / Journal of Combinatorial Theory, Series B 92 (2004) 325--357}
\noindent be a hypergraph and $\theta \ge 1$ an integer. A tangle of order $\theta$ in $G$ is a set $\mathcal{T}$ of separations of $G$, each of order $< \theta$, such that

\hspace*{12pt}$\bullet$ for every separation $(A, B)$ of $G$ of order $< \theta$, $\mathcal{T}$ contains one of $(A, B)$, $(B, A)$,

\hspace*{12pt}$\bullet$ if $(A_i, B_i) \in \mathcal{T}$ $(i = 1, 2, 3)$ then $A_1 \cup A_2 \cup A_3 \ne G$,

\hspace*{12pt}$\bullet$ if $(A, B) \in \mathcal{T}$ then $V(A) \ne V(G)$.

Let us mention one lemma that we shall need later.

\textbf{2.1.} \textit{Let $G$ be a hypergraph, let $G' \subseteq G$ and let $\mathcal{T}'$ be a tangle in $G'$ of order $\theta$. Let $\mathcal{T}$ be the set of all separations $(A, B)$ of $G$ of order $< \theta$ such that $(A \cap G', B \cap G') \in \mathcal{T}'$. Then $\mathcal{T}$ is a tangle in $G$ of order $\theta$.}

The proof is clear.

A \textit{tie-breaker} in a hypergraph $G$ is a function $\lambda$ which maps each separation $(A, B)$ of $G$ to some member $\lambda(A, B)$ of a linearly ordered set $(\Lambda, \le)$ (we call $\lambda(A, B)$ the \textit{$\lambda$-order} of $(A, B)$) in such a way that for all separations $(A, B)$, $(C, D)$ of $G$,

\hspace*{12pt}$\bullet$ $\lambda(A, B) = \lambda(C, D)$ if and only if $(A, B) = (C, D)$ or $(A, B) = (D, C)$,

\hspace*{12pt}$\bullet$ either $\lambda(A \cup C, B \cap D) \le \lambda(A, B)$ or $\lambda(A \cap C, B \cup D) < \lambda(C, D)$,

\hspace*{12pt}$\bullet$ if $|V(A \cap B)| < |V(C \cap D)|$ then $\lambda(A, B) < \lambda(C, D)$.

Let $\lambda$ be a tie-breaker in a hypergraph $G$. If $\mathcal{T}_1, \mathcal{T}_2$ are tangles in $G$ with $\mathcal{T}_1 \nsubseteq \mathcal{T}_2$ and $\mathcal{T}_2 \nsubseteq \mathcal{T}_1$, then there is a unique $(A, B) \in \mathcal{T}_1$ such that $(B, A) \in \mathcal{T}_2$ of minimum $\lambda$-order, called the \textit{$(\mathcal{T}_1, \mathcal{T}_2)$-distinction}.

A \textit{march} in a set $V$ is a finite sequence of distinct elements of $V$; and if $\pi$ is the march $v_1, \ldots, v_k$, we denote the set $\{v_1, \ldots, v_k\}$ by $\bar{\pi}$. We denote the null march by $0$. A \textit{rooted hypergraph} $G$ is a pair $(G^-, \pi(G))$ where $G^-$ is a hypergraph and $\pi(G)$ is a march in $V(G^-)$. We define $V(G) = V(G^-)$, $E(G) = E(G^-)$. If $G$ is a rooted hypergraph, a \textit{tangle} in $G$ is a tangle in $G^-$, and a \textit{tie-breaker} in $G$ is a tie-breaker in $G^-$.

A \textit{separation} of a rooted hypergraph $G$ is a pair $(A, B)$ of rooted hypergraphs such that $(A^-, B^-)$ is a separation of $G^-$, $\bar{\pi}(A) = V(A \cap B)$, and $\pi(B) = \pi(G)$. If $G, A$ are rooted hypergraphs, we write $A \subseteq G$ if $A^- \subseteq G^-$. If $A \subseteq G$, we say $A$ is \textit{complemented} if there exists $B \subseteq G$ such that $(A, B)$ is a separation of $G$, and we define $G \setminus A = B$. A \textit{rooted location} in a rooted hypergraph $G$ is a set $\mathcal{L}$ of complemented rooted hypergraphs $A$ with $A \subseteq G$ such that $E(A_1^- \cap A_2^-) = \emptyset$ and $V(A_1^- \cap A_2^-) = \bar{\pi}(A_1) \cap \bar{\pi}(A_2)$ for all distinct $A_1, A_2 \in \mathcal{L}$. Its \textit{order} is $\max(|\bar{\pi}(A)| : A \in \mathcal{L})$, or $0$ if $\mathcal{L} = \emptyset$. If $\mathcal{L}$ is a rooted location in $G$, we define $\mathcal{L}^- = \{(A^-, (G \setminus A)^-) : A \in \mathcal{L}\}$, and we define $M(G, \mathcal{L})$ to be $\cap((G \setminus A)^- : A \in \mathcal{L})$ if $\mathcal{L} \ne \emptyset$, and to be $G^-$ if $\mathcal{L} = \emptyset$.

Let $G$ be a rooted hypergraph, let $\mathcal{T}$ be a tangle in $G$, and let $\lambda$ be a tie-breaker in $G$. A rooted location $\mathcal{L}$ in $G$ is said to \textit{$\theta$-isolate} $\mathcal{T}$ if $\theta \ge 1$, $\mathcal{L}$ has order $< \theta$, $\mathcal{L}^- \subseteq \mathcal{T}$, and for each $A \in \mathcal{L}$, and for every tangle $\mathcal{T}'$ in $G$ of order $\ge \theta$ with $((G \setminus A)^-, A^-) \in \mathcal{T}'$, the $(\mathcal{T}, \mathcal{T}')$-distinction $(C, D)$ satisfies $C \subseteq A^-$ and $(G \setminus A)^- \subseteq D$.

\sectionhead{3. Patchworks}

If $V$ is a finite set we denote by $K_V$ the complete graph on $V$, that is, the simple graph with vertex set $V$ and edge set the set of all subsets of $V$ of cardinality $2$, with the natural incidence
